\documentclass[10pt,a4paper]{article}
\usepackage[T1]{fontenc}
\usepackage{lmodern}
\usepackage{textcomp}
\usepackage[margin=2.2cm]{geometry}
\usepackage{graphicx}
\graphicspath{{figures/}{./}}
\usepackage{amsmath}
\usepackage{caption}
\usepackage{enumitem}
\usepackage{hyperref}
\captionsetup{font=small}
\setlength{\textfloatsep}{10pt}
\setlength{\intextsep}{8pt}

\begin{document}

\begin{titlepage}
  \centering
  \vspace*{4cm}
  {\Huge\bfseries Which Dutch supermarket\\[4pt] is the cheapest?\par}
  \vspace{0.8cm}
  {\Large Comparing the six largest chains,\\ November 2022 -- October 2026\par}
  \vspace{2.5cm}
  {\large Jaimy Hereygers\par}
  \vspace{0.2cm}
  {\large I6285837\par}
  \vspace{0.6cm}
  {\normalsize Code: \url{https://github.com/jhereygersstudent/Mat2007}\par}
  \vfill
  {\large MAT2007 Introduction to Programming\par}
  \vspace{0.2cm}
  {\large Maastricht Science Programme, Maastricht University\par}
  \vspace{0.6cm}
  {\large October 2026\par}
\end{titlepage}

\setcounter{page}{1}

\section*{Introduction}
Groceries in the Netherlands have become noticeably more expensive in recent years,
while almost every supermarket advertises itself as the cheapest option. Discounters
such as Dirk, Lidl and Aldi build their image on low prices, while Albert Heijn and
Jumbo promise to match them with special offers.

This report addresses a simple question: \textbf{Which of the six largest Dutch
supermarket chains offered the cheapest everyday shopping basket between November 2022
and October 2026, and did the ranking change over time?}

To answer this, public price data from the price-comparison site
\emph{checkjebon.nl}\footnote{\url{https://www.checkjebon.nl}} are analysed in R; the code is available on GitHub~\cite{code}.

\section*{Data sample}
The dataset contains the daily online product lists of Dutch supermarkets, archived on
GitHub~\cite{cjb}. The first snapshot of every month from November 2022 to October 2026
was used (48 snapshots, about 3.1 million prices), for Albert Heijn, Jumbo, Lidl, Plus,
Aldi and Dirk. Lidl only appears from December 2025. The relevant fields are:
\begin{itemize}[itemsep=0pt,topsep=2pt]
  \item \texttt{n}: product name
  \item \texttt{p}: price in euro
  \item \texttt{s}: pack size (e.g.\ \texttt{1 l}, \texttt{500 g})
\end{itemize}
The data required cleaning: pack sizes were converted to grams, millilitres or pieces,
and for each of 30 everyday products (milk, eggs, cheese, pasta, oats, coffee, \dots)
the cheapest matching product per chain and month was selected. Each price was divided
by the median price of that item across the six chains, and the mean of these ratios
gives the chain's price level. Months without data repeat the previous month.

\begin{figure}[h]
  \centering
  \includegraphics[width=0.9\linewidth]{plot1_basket_over_time.png}
  \caption{Cost of the 30-item basket per supermarket chain and month,
  November 2022 -- October 2026.}
  \label{fig:time}
\end{figure}

\section*{Results and discussion}
Figure~\ref{fig:time} shows the basket cost of each chain over time. Dirk sat lowest
overall and came out on top in 30 of the 48 months, with a basket about 7.5\% under the
average of the group. Lidl looked almost the same when it was there, but the numbers
only ran for eleven months. Aldi was a bit behind them. Jumbo and Albert Heijn were only
slightly over the average, and Plus was always the highest by quite a bit:
\begin{equation*}
\Delta_{\text{Dirk}} = -7.5\%, \qquad \Delta_{\text{Lidl}} = -6.6\%, \qquad
\Delta_{\text{Aldi}} = -3.3\%, \qquad \Delta_{\text{Plus}} = +9.3\%
\end{equation*}

Over time the gap got smaller. Early on Dirk was way ahead, at 19\% below average in
late 2022, and later the difference shrank to around 4\%. In the last year Lidl and
Jumbo traded places near the bottom more often. Figure~\ref{fig:rank} shows the spread
of these monthly differences per chain.

\begin{figure}[h]
  \centering
  \includegraphics[width=0.75\linewidth]{plot2_ranking.png}
  \caption{Monthly difference of each chain's basket to the average of the six chains.
  Grey dots are single months, the large dot is the median.}
  \label{fig:rank}
\end{figure}

It seems the bigger chains started closing in on the discounters, which makes sense if
they were trying to keep customers. Some parts of the data stayed flat for long
stretches because prices were not updated, which probably made the 2026 numbers for
Lidl and Jumbo look a little better than they really were. Products were matched by
name, and offers and loyalty prices were left out, so the real cost for shoppers could
be different. Still, the overall order stayed fairly steady, with the three discounters
under the others and Plus at the top every single month.

\begin{thebibliography}{9}
\bibitem{cjb} Checkjebon.nl. \emph{Supermarket price data}. Available at:
\url{https://github.com/supermarkt/checkjebon} (accessed October 2026).
\bibitem{code} J.~Hereygers. \emph{MAT2007 final project: R code and README}. Available at:
\url{https://github.com/jhereygersstudent/Mat2007}.
\end{thebibliography}

\end{document}
